% =====================================================================
\chapter{Introduction}
% =====================================================================

\section{The questions}

The following three questions share one structure. They are stated here in the form in
which they can be discussed, proved and falsified.

\begin{problem}[Generality of low-level determinism]\label{prob:low}
Software at the execution level is strikingly uniform: ABIs and calling conventions, the
call stack, heap and object lifetime, control and data flow, synchronisation primitives
and thread scheduling. That uniformity is imposed by a shared computational model,
hardware architecture, language runtime and operating-system interface. Does the same
determinism extend upwards? Do most large software systems, normally biased and normally
engineered, converge on the same or similar abstractions, structures and architectures?
\end{problem}

\begin{problem}[Coercion by scale]\label{prob:scale}
When a system crosses some order of magnitude, does it necessarily begin to design and
optimise for recurrent problems of one kind, and thereby grow a unified internal
infrastructure? An operating-system kernel defines how memory is managed and allocated,
how threads and processes are defined and scheduled, and what runtime interface is
offered upwards. A browser, a game engine, an asynchronous runtime, a distributed store
and a relational engine address the same class of questions at different levels. Once the
code base is large enough, do all of them exhibit the same kind of reasoning---unified
function, unified interfaces, managed planning, runtime structure?
\end{problem}

\begin{problem}[Comprehension and design]\label{prob:mind}
Software is written by humans and machines, but it is not easy to understand. A large
system can hardly be read sequentially from the micro level. At what level of abstraction
should it be approached, and how does comprehension lead into development without
dissolving into a sea of code? The question bears equally on machines: an intelligence
cannot descend into unbounded low-level information, and must retain an understanding of
first principles---call this \Lk{0}---together with the levels above it, \Lk{1}, \Lk{2},
\Lk{3}, and so on. What an intelligence needs is compressed information. Conversely, how
should a large complex system be designed so that its purpose is achieved?
\end{problem}

The third question has a contemporary variant that moves it from comprehension to
warranty:

\begin{problem}[Correctness without review]\label{prob:trust}
If machine generation drives the cost of code towards zero and no human reads the
implementation line by line, how is correctness guaranteed \emph{mechanically} and
\emph{non-self-referentially}? How is it guaranteed that the stated purpose is achieved,
and that no unauthorised side effect accompanies it? A concrete instance: a hand-written
call takes 1\,ns; a generated implementation takes 1\,$\mu$s. Functionally identical,
purpose achieved, and a performance side effect introduced.
\end{problem}

\section{Three gaps in existing treatments}

Each question has local counterparts in the literature: architectural convergence in
pattern languages and complex-systems theory, coercion by scale in information hiding and
near-decomposability, comprehension and design in cognitive models and multi-view
modelling, correctness in program logic, refinement calculi and verified compilation.
Treating them as a whole exposes three structural gaps.

\paragraph{The strength of the convergence claim is not graded.}
``Large software architectures converge'' contains at least three propositions of very
different strength: recurrent local architectural patterns (well supported by engineering
experience); measurable structural similarity above a random baseline after abstraction
(plausible and worth systematic testing); and convergence of nearly all large systems on
one high-level architecture (unsupported, and probably too strong). Conflating them makes
the claim untestable.

\paragraph{Necessity claims lack a load-bearing structure.}
``Scale forces a unified infrastructure'' is a necessity claim, but it is usually stated
as an empirical observation, and therefore cannot answer where the necessity comes from,
what fixes the threshold, or why some equally large systems grow no kernel. The answers
given below are that the necessity comes from the completeness of the operational
semantics, the threshold from an amortisation inequality, and the apparent exceptions
from systems that \emph{adopt} a kernel rather than build one.

\paragraph{Comprehension, design and assurance are three readings of one structure.}
Comprehension is treated as a cognitive problem, design as an engineering problem and
assurance as a logical problem, in three separate literatures. They are in fact different
reading directions of one object: a level is the unit into which assurance obligations
decompose; comprehension recovers that decomposition and design prescribes it.

\section{Core theses}

Four theses organise the paper. Each is annotated with its formal location and its
strength: \textbf{[P]} proved, \textbf{[A]} argued under stated assumptions,
\textbf{[H]} advanced as a falsifiable hypothesis.

\begin{principle}[Compression]\label{th:compression}\textbf{[A]}
A complex system is comprehensible to a bounded intelligence if and only if it is
compressible. Comprehension is task-constrained lossy compression: description length is
minimised subject to preserved task-relevant predictive power. The optimal descent depth
is therefore the stationarity condition of a rate--distortion Lagrangian, that is, the
point at which marginal predictive gain vanishes.
\end{principle}

\begin{principle}[Necessity]\label{th:necessity}\textbf{[A]}
Any system with concurrently live entities and contended finite resources, whose
operational semantics is well defined, has already defined six functions: identity and
lifecycle, allocation and reclamation, scheduling, naming and addressing, communication
and synchronisation, and fault isolation and recovery. These may be built or adopted from
a host---an operating system, a language runtime, a framework---but they cannot be omitted
from the semantics.
\end{principle}

\begin{principle}[Separation]\label{th:separation}\textbf{[A]}
Generation may be free, probabilistic, open and untrusted; admission must be
deterministic, mechanical and exhaustively inspectable. Guarantees on the formal layer
(specification $\to$ implementation $\to$ binary $\to$ execution) come from mechanical
decision by a minimal trusted kernel. Guarantees on the intent layer (intent $\to$
specification) come from adversarial consensus, and that layer admits no mechanical
solution. The two are stitched by artefact binding: only the verified binary is admitted
to run, and the environmental assumptions are checked or monitored on it.
\end{principle}

\begin{principle}[Design]\label{th:design}\textbf{[A]}
Architecture is manufactured separation of scales. Nature obtains levels free from scale
separation and renormalisation; software levels do not emerge spontaneously and must be
paid for in interfaces, encapsulation and invariants, and they leak. Design is therefore
forward compression and comprehension is inverse compression, two reading directions of
one structure. The measurable objective of architecture is the expected read-set
$\mathbb{E}[\rho]$ of a typical change.
\end{principle}

The four theses and their formal results are referred to collectively as the
\term{compression--assurance framework}. The name is taken from the two ends of the claim:
\emph{compression} is the mechanism of comprehension and design, \emph{assurance} is the
mode in which trust is anchored. The central assertion is that a complex system is
comprehensible to a bounded agent if and only if it is compressible, and that it is
compressible if and only if it admits a separation of scales in which every level has an
explicitly declared oblivious set.

\section{Epistemic discipline}

The material crosses engineering intuition, complex-systems theory, information theory,
mathematical logic and social epistemology, and the risk of overreach is high. Five rules
are observed throughout.

\begin{enumerate}[label=(\roman*)]
\item \term{Strength labelling.} Every claim is labelled as a theorem, proposition,
      conjecture, hypothesis or heuristic. Heuristics are not written as theorems.
\item \term{Non-claims list.} Chapter \ref{ch:objections} states explicitly what the paper
      does \emph{not} claim, including several ``authoritative theorems'' that are declined.
\item \term{Explicit assumptions.} Any conclusion that depends on an environmental
      assumption states it, and any guarantee that has a boundary marks it.
\item \term{Falsifiability.} Chapter \ref{ch:program} reduces the core claims to fourteen
      falsifiable hypotheses with operationalisations, controls, falsifiers and risks.
\item \term{Verification not isomorphic to the conclusion.} Where the paper argues that
      generation must be untrusted and verification independent, that rule applies to the
      paper itself: its load-bearing propositions must be independently checkable and must
      not depend on any unproved conjecture.
\end{enumerate}

Rule (v) constrains the paper's own rhetoric. Several claims have been removed from the
load-bearing structure: G\"odel incompleteness and Rice's theorem are not used to prove
that a gap between intent and specification is principled
(Chapter \ref{ch:intent}, \S\ref{sec:intent-gap}); the open status of $\mathrm{P}$ versus
$\mathrm{NP}$ is not used to support ``verification is cheaper than generation''
(\S\ref{sec:asymmetry}); and near-decomposability is not treated as a universal law
(Chapters \ref{ch:convergence} and \ref{ch:objections}). Each of these gains strength by
being unloaded rather than reinforced.

\section{Contributions}

Formal contributions:

\begin{enumerate}[label=(C\arabic*)]
\item The \emph{oblivious set} of a level as the kernel of its abstraction map, which makes
      ``what this level need not care about'' a computable object
      (Definition \ref{def:oblivious}).
\item Galois connections as the condition for inter-level semantic preservation
      (Definitions \ref{def:galois}, \ref{def:soundness}, Proposition \ref{prop:soundness}).
\item A \emph{contract operad} whose well-definedness reduces the composability of local
      guarantees to well-typed colouring and acyclicity of the interface graph
      (Definition \ref{def:operad}, Proposition \ref{prop:operad-compose}).
\item A \emph{level family} $\pi:\mathcal{L}\to\mathcal{P}$ that separates cross-domain
      isomorphism of the six functions from transport along problem shape
      (Definition \ref{def:fibration}, Proposition \ref{prop:transport}).
\item A \emph{coarse-graining criterion} under which the convergence thesis can be
      falsified inside a single project (Definition \ref{def:coarsegrain},
      Conjecture \ref{conj:relevant}).
\item An amortisation inequality comparing $\Theta(nm)$ direct interactions with
      $\mathcal{O}(nk+m)$ substrate interactions, giving the break-even threshold of internal
      infrastructuralisation (Theorem \ref{th:integration}).
\item Conditional necessity of the six functions from the well-definedness of the
      operational semantics (Theorem \ref{th:six}), and the resulting build-or-adopt
      criterion (Corollary \ref{cor:buildadopt}).
\item A criterion for over- and under-constraint that renders ``unique below, manifold
      above'' as solution-set geometry (Proposition \ref{prop:overconstrain}).
\item Expected read-set as the measurable objective of architecture, with a bound showing
      that invariant centralisation lowers it, stated relative to the enforcement sites that
      no independent evidence settles (Definition \ref{def:readset}, Definition
      \ref{def:enforce}, Theorem \ref{th:centralise}).
\item Technical debt as a description-length gap (Definition \ref{def:debt}), with a
      rate--distortion reading of refactoring.
\item An identifiability criterion for whether comprehension can converge, separating
      insufficient observation from insufficient sampling, of which the linear observability
      rank is the finite-dimensional specialisation (Definition \ref{def:observability},
      Proposition \ref{prop:observability}).
\item A governable-complexity stability condition and a scaling law for required
      headcount (Definition \ref{def:governable}, Proposition \ref{prop:exponent}).
\item A Hoeffding bound for consensus reliability and the proof that a correlated error
      mode dominates it, reducing the governing parameter to an incentive-design problem
      (Theorem \ref{th:condorcet}, Proposition \ref{th:correlated}, \S\ref{sec:mechanism}).
\item The observation that the mechanical and consensus layers are two enrichments of one
      structure, and that no information-preserving change of base runs from the second to the
      first; consequently reliability must be \emph{budgeted} across the stages that must all hold
      rather than gated stage by stage, a per-stage threshold giving a composite bound that grows
      linearly and is vacuous once the number of stages reaches the reciprocal of the
      threshold (\S\ref{sec:enrichment}, Principle \ref{pr:budget}). The mathematical content here
      is Lawvere's, not ours; the contribution is its application to the two-layer structure and
      the budget rule that follows.
\item A consolidated assurance condition combining two-layer trust with artefact binding
      (Definition \ref{def:assurance}, Proposition \ref{prop:assurance}).
\end{enumerate}

Methodological contributions:

\begin{enumerate}[label=(M\arabic*)]
\item A descent protocol that converts comprehension of a large system from sequential
      reading into layered modelling and active identification.
\item An executable list of seam detectors for locating real module boundaries in an
      unfamiliar code base.
\item A design procedure that replaces adjectival goals with an optimisable objective and
      a stability condition.
\item Fourteen falsifiable hypotheses with experimental designs.
\item A cross-disciplinary research position that couples social epistemology to
      proof-carrying code, which the literature map identifies as unoccupied.
\end{enumerate}

\section{Structure}

Chapter \ref{ch:framework} establishes the formal language: levels, abstraction maps,
oblivious sets, contracts, traces, assurance and trust. Chapter \ref{ch:convergence}
addresses Problem \ref{prob:low}, grading the convergence claim and stating its strength
law. Chapter \ref{ch:infrastructure} addresses Problem \ref{prob:scale}, proving the
conditional necessity of the six functions and deriving the break-even threshold at which a substrate pays for itself.
Chapter \ref{ch:comprehension} addresses the first half of Problem \ref{prob:mind}:
comprehension as inverse compression. Chapter \ref{ch:design} addresses the second half:
design as forward compression. Chapter \ref{ch:intent} addresses Problem \ref{prob:trust}
with a verifiable refinement chain and a two-layer trust structure. Chapter
\ref{ch:levels} promotes levels from a comprehension aid to the unit into which assurance
obligations decompose. Chapter \ref{ch:program} states the falsifiable programme. Chapter
\ref{ch:literature} positions the work. Chapter \ref{ch:objections} lists objections and
limits. Chapter \ref{ch:conclusion} concludes.

\section{Terminology}

Several terms drift between communities. They are fixed once here; the full glossary is
Appendix \ref{app:glossary}, and the Chinese rendering is fixed term by term in the
companion terminology file.

\begin{longtable}{L{3.1cm}L{11.4cm}}
\toprule
\textbf{Term} & \textbf{Meaning in this paper} \\
\midrule
\endhead
level & A tuple (representation space, behaviour set, observation map, environment
  assumptions, cost model, invariants), defined relative to a study object and a task;
  a cognitive scale, not a physical directory. \\
oblivious set & The differences a level can never distinguish; the kernel of its
  abstraction map. \\
separation of scales & The condition under which adjacent levels each support a valid
  description and the lower level's implementation details do not appear in the upper
  level's semantics. \\
internal infrastructuralisation & The process by which a system promotes recurrent
  management problems to unified mechanisms on which the upper layers then depend. \\
the six functions & Identity and lifecycle; allocation and reclamation; scheduling;
  naming and addressing; communication and synchronisation; fault isolation and recovery. \\
read-set & The set of modules that must be read to complete a typical change. \\
assurance & A traceable evidential relation holding a property within stated premises and
  a stated scope. \\
mechanical assurance & Assurance by a single deterministic, exhaustively inspectable
  decision procedure. \\
consensus assurance & Assurance by multiple independent, adversarial participants. \\
artefact binding & The identity relation between the verified artefact and the executed
  artefact. \\
intent gap & The distortion between an intent and its formal projection. \\
\bottomrule
\end{longtable}

One terminological caution is load-bearing. A \emph{level} is a cognitive scale, not a
physical tier. The same component occupies different levels in different tasks: an
asynchronous runtime is infrastructure to an application developer and a complex system in
its own right to its maintainer. Levels are defined relative to the study object and the
task (Chapter \ref{ch:levels}).
